\backhead{FURTHER EXERCISES}{khred}
\begin{exlist}
\item A field of 6 ha of the sandy loam ($\alpha=0.075$ cm$^{-1}$, $n=1.89$, $\theta_r=0.065$, $\theta_s=0.41$) has a root zone of 0.7 m. (a) Find $\theta$ at field capacity (100 cm) and at the wilting point (15\,849 cm). (b) Find $\TAW$ in mm. (c) With $p=0.5$ and $ET_c=5.5$ mm\,d$^{-1}$, find the interval and the volume of each irrigation at an efficiency of 0.8.\exans{(a) 0.122, 0.066; (b) 39 mm; (c) 3.6 d, 1475 m$^3$}
\item Tensiometers at 30 and 60 cm read suctions of 250 and 120 cm (datum at the surface). (a) Find the total heads. (b) In which direction does water flow? (c) If $K=0.05$ cm\,d$^{-1}$ in this range, find the flux.\exans{(a) $-280$ and $-180$ cm; (b) upward; (c) 0.17 cm\,d$^{-1}$ upward}
\item For the loam ($\alpha=0.036$, $n=1.56$) find the suction at which $S_e=0.5$, and the pore radius that is just full there.\exans{87 cm; 17.1 $\mu$m}
\item A rainfed year on a loam gives $P=640$, $T=290$, $E_s=170$, $E_o=55$, $D=110$ and $Q=8$ mm. (a) Find $\Delta S$. (b) An irrigation of 150 mm is added; the extra water is 40 per cent transpired, 35 per cent evaporated and 25 per cent drained. Find the new $T$, $E_s+E_o$ and $D$.\exans{(a) $+7$ mm; (b) 350, 277.5, 147.5 mm}
\item By how much does $K_r$ of the sand ($\alpha=0.145$, $n=2.68$) fall between 50 and 300 cm? Compare with the loam over the same range.\exans{sand: factor $66\,132$; loam: factor $271$}
\end{exlist}
\backhead{ESSAY QUESTIONS}{khred}
{\small\itshape Answer each question in about one page, with equations, a sketch and at least one number from this chapter.}\par\vspace{3pt}
\begin{exlist}
\item Field capacity and wilting point are conventions, not constants of the soil. Discuss what each convention assumes and how the choice of the suctions changes the computed available water.
\item Explain the components of the soil-water potential and show which of them drive flow in the soil, which drive uptake by roots, and which matter in a saline soil.
\item Compare three methods of measuring soil water (gravimetric, TDR, tensiometer). For each, state what is measured, what is inferred, and the main sources of error.
\item Discuss the meaning of the van Genuchten parameters $\alpha$ and $n$ in terms of the pores, and show how each changes the available water and the conductivity of a soil.
\item Why is the conductivity of an unsaturated soil not a constant? Explain the consequences for irrigation of a sand and of a clay.
\item Write the water balance of a root zone and discuss which of its terms are measured, which are modelled and which are obtained as a residual. What are the dangers of the residual?
\end{exlist}
